\documentclass[a4paper,10pt]{article}
%\documentclass[a4paper,10pt]{scrartcl}

\usepackage[utf8]{inputenc}
\usepackage{hyperref}
\usepackage{graphicx}
\usepackage{textcomp}
\usepackage{amsmath}

\hypersetup{
    colorlinks,
    citecolor=black,
    filecolor=black,
    linkcolor=black,
    urlcolor=blue
}

\title{Dokumentace projektu do předmětu \\* Modelování a simulace\\* zadání 8: drůbežárna}
\author{Jiří Zahradník a Veronika Lysáková}
\date{}

\pdfinfo{%
  /Title    ()
  /Author   ()
  /Creator  ()
  /Producer ()
  /Subject  ()
  /Keywords ()
}

\begin{document}

\maketitle
\includegraphics[scale=0.5]{fitnewb.png}
\pagebreak

\renewcommand{\contentsname}{Obsah}
\tableofcontents

\renewcommand{\figurename}{Obrázek}


\pagebreak

%
% úvod dokumentace
%

\section{Úvod}
V tomto textu bude čtenář seznámen s vytvářením modelu pro diskrétní simulaci 
života broilerů a ekonomické udržitelnosti drůbežárny.
\par 
Je zde popsáno studium autorů, kde autoři získávají potřebné informace
pro vytvoření simulačního modelu na základě reálného modelu drůbežárny.
\par
Na základě experimentů bude demonstrována ekonomická udržitelnost
drůbežárny za současných ekonomických podmínek.

\subsection{Zdroje informací pro modelování\cite{modelovani}}\label{zdroje}
Jako zdroj informací autoři využili oficiální materiály firmy Aviagen\cite{aviagen}.
Jedná se o firmu zajišťující dovoz jednodenních kuřat do drůbežáren ve 130 zemích světa.
Autoři se zaměřili konkrétně na broilery značky Ross\textsuperscript{\textregistered},
jejichž podmínky pro pěstování jsou popsány v dokumentu poskytnutém firmou Aviagen\cite{ross}
na jejich oficiálních webových stránkách. Autoři následně využili tyto informace k
vytvoření abstraktního modelu\cite{abstract_model}, který byl posléze převeden na simulační model\cite{simulation_model}
simulující farmu pro výkrm broilerů Ross\textsuperscript{\textregistered}. 

\subsection{Validace a verifikace}\label{validaceVerifikace}
TODO model zvalidovan pomoci smernic EU


%
% Téma a použité metody
%

\section{Rozbor tématu, použitých metod a technologií}\label{rozbor}
Drůbežáren může být několik druhů. Patří mezi ně drůbežárny soustředící 
na produkci nosných slepic, drůbežárny na produkci vajec, drůbežárny 
na produkci kuřat určených k rychlému výkrmu, tzv. broilerů a další.
V této modelové studii se autoři zaměřili na drůbežárny pro produkci broilerů.
\par
Kuřata určena pro výkrm, dovezena do drůbežárny jsou stará jeden den 
a važí 40 gramů. Prvních 10 dní jsou krmena krmivem typu Starter. Toto krmivo 
musí zajistit dostatečný obsah živin pro kuřata, která v prvních deseti dnech
po vylíhnutí konzumují nejmenší objem krmiva, avšak jejich nároky na výživnost krmiva
jsou největší. Tento fenomén je způsoben přizpůsobování kuřete k životu mimo skořápku.
V tomto stádiu se rozhoduje, zda kuře bude mít dostatečnou chuť k jídlu, zda bude zdravé
a zda bude správně růst.
\par
Druhá část výkrmu je dlouhá 14 dní, od 11. dne po 25. den. V této části se využívá krmivo
Grower, a je podáváno ve formě větších granulí, než tomu bylo u krmiva Starter.
Rapidně se zvyšuje váhový přírustek broilerů, jelikož obsahuje menší procento bílkovin
a aminokyselin, avšak jeho metabolizovatelná nergetická hodnota se zvýšila. Při přechodu
mezi krmivy Starter a Grower je třeba pozorně sledovat kuřata a musí se předejít zmenšení
příjmu potravy nebo oslabení růstu.
\par
Poslední část výkrmu nastává okolo 25. dne věku kuřat. Krmivo Grower nahrazuje krmivo Finisher.
Toto krmivo tvoří největší část příjmu potravy kuřat a tím největší část nákladů na výkrm kuřat.
Proto musí být krmivo Finisher optimalizované na finanční návratnost. Ke konci výkrmu
je třeba vysadit farmaceutické aditiva ať se maso broilerů pročistí 
a nezůstávají v něm rezidua. Avšak je třeba zachovat energetickou hodnotu krmiva z důvodu
udržení hmotnostního přírůstku kuřat. Kuřata jsou po čtyřiceti až dvaačtyřiceti dnech zabita
a jejich maso je zpracováno.

\subsection{Postupy}\label{postupy}
Po nastudování a sepsání potřebných informací autoři vytvořili abstraktní model v podobě
Petriho sítě\cite{petriho_sit} života kuřete. Viz. obrázek \ref{obr:petri_kure}.
\par
Pro vytvoření simulačního modelu měli autoři předepsané jazyky C nebo C++. Autoři se rozhodli
pro C++ kvůli možnosti objektového návrhu, který umožňuje provést abstrakci pomocí objektů
a tato abstrakce umožňuje objektu specifické chování na základě vztupních podmínek.
Při porovnání s jazykem C, kde je možno implementovat pseudoobjektový přístup, je objektový 
přístup jazyka C++ intuitivnější a umožňuje jednodušší práci s jednotlivými objekty.
Každé kuře je reprezentováno objektem a s těmito objekty pracuje objekt továrna, který uchovává
data o chovu kuřat. Více o konceptu v kapitole \ref{koncepce}.

\subsection{Původy metod}\label{puvod}
Informace pro práci s jazykem C++ autoři čerpali ze semináře jazyka C++\cite{cpp}
přednášeného doktorem Peringerem. Informace pro práci s objekty byly čerpány
z předmětu IPP\cite{ipp}.





\section{Koncepce}\label{koncepce}
TODO
Jelikož nás pouze zajímá ekonomická situace drůbežárny, můžeme z modelu vyloučit detaily
o růstu kuřete. Pro vyhodnocení ekonomické situace nám stačí pouze množství snězeného krmiva,
spotřebované vody, výsledná váha kuřete, kupní cena jednodenního broileru Ross\textsuperscript\textregistered
a velkoobchodní cena vykrmeného broileru.
\par
Koncepční model\cite{concept_model} je vytvořen jako diskrétní model, kde časovou osu reprezentují dny věku kuřete.
Kuře v modelu je reprezentováno jako jako Markovův proces\cite{markovuv_proces},
kdy každý den proces mění svůj stav na základě svého aktuálního věku.
\par
Život kuřete je modelován Petriho sítí, která zobrazena je na obrázku \ref{obr:petri_kure}.
Po dovezení kuřete začíná dvaačtyřiceti denní výkrm, kde je kuře postupně krmeno, napájeno, v jehož důsledku
přibírá na váze. Kuře taky může umřít před datem zpracování a to je znázorněno přechodem\cite{prechod} s pravděpodobností 0.05\%.
Pokud přežije, přechází do dalšího dne a cyklus se opakuje. Jak bylo již řečeno, časová osa modelu je reprezentována dny
věku kuřat. O posuny vpřed na časové ose se stará drůbežárna, jejichž lze vidět na obrázku \ref{obr:drubezarna}. 
Po zpracování jsou data o spotřebě kuřete předána drůbežárně, ta je vyhodnotí a na tomto základě vytvoří finanční zrávu.

\subsection{Zdůvodnění způsobu vyjádření konceptuálního\newline modelu}
Na obrázku \ref{obr:petri_kure} je zobrazen model života kuřete. Petriho síť byla vybrána, jelikož umožňuje
jednoduché časování přechodů a síť lze bez větších obtíží transformovat na simulační model v jazyce C++, kde přechody
označují volání metod odpovídajícího objektu v simulačním modelu.

\subsection{Konceptuální model}
\begin{figure}[h]
\caption{Petriho síť modelující život kuřete}
\includegraphics[scale=0.4]{kure.png}
\label{obr:petri_kure}
\end{figure}

Na obrázku \ref{obr:drubezarna} můžeme vidět model drůbežárny samotné v podobě Petriho sítě. Můžeme vidět, že po přijeťí kuřat
se model dostane do stavu, kde se rozhoduje, jestli se bude posouvat po časové ose anebo jestli zpracuje
kuřata. Proměnná A vyjadřuje počet tokenů(teček), které reprezentují množinu kuřat vynásobenou počtem dní výkrmu,
nebo-li: 
$A=\textit{počet kuřat }*\textit{ dny výkrmu}$.
Dokud model iteruje, volá se metoda simulačního modelu \textit{nextDay()}, které posouvá simulační model po časové ose doprava
a až bude proměnná A rovna nule, provede se přechod, který symbolizuje volání metody, která zpracuje kuřata. Následně 
se vytvoří finanční zpráva a simulace se ukončí.


\begin{figure}[h]
\caption{Petriho síť modelující provoz drůbežárny}
\includegraphics[scale=0.4]{drubezarna.png}
\label{obr:drubezarna}
\end{figure}

\pagebreak

\section{Architektura simulačního modelu}\label{architektura}
Architektura simulačního modelu je postavena na objektovém přístupu. Simulační model se skládá ze dvou tříd,
třídy factory, a třídy chicken. Z třídy \textit{factory} se vytvoří objekt drůbežárny, a z třídy \textit{chicken}
se vytvoří jednotlivá kuřata, která se seřadí do vektoru \textit{chick}. S tímto vektorem poté pracuje továrna tak,
že nad ním volá metodu \textit{nextDay()}(lze vidět na obrázku \ref{obr:drubezarna}), která zašle objektu kuřete
zprávu o příchodu následujícího dne. Po přijetí této zprávy objekt kuře simuluje krmení, pití a tloustnutí.
Poté objekt kuře zkontroluje booleovskou hodnotu signalizující, zda je nakaženo a má umřít anebo žije další den.
Tato booleovská hodnota je nastavena na true v 0.05\% případů. Kontrola se provádí každý den a dohromady to dává ~2\% celkový úhyn.
Pokud kuře uhyne, zvýší se čitač v objektu drůbežárny a s mrtvým objektem kuře se nadále nepracuje.
Po ukončení výkrmu se zavolá metoda drůbežárny \textit{statistics()}, která vyhodnotí údaje o výkrmu a vytvoří vinanční analýzu.

\section{Podstata simulačních experimentů a jejich průběh}\label{experimenty}
TODO

\section{Shrnutí simulačních experimentů a závěr}\label{zaver}
TODO

\renewcommand\listfigurename{Seznam obrázků}
\listoffigures
\addcontentsline{toc}{section}{\listfigurename}


\renewcommand{\refname}{Zdroje}
\bibliographystyle{unsrt}
\bibliography{dokumentace}
\addcontentsline{toc}{section}{\refname}


\end{document}

\begin{thebibliography}{99}

\bibitem{markovuv_proces}
\bibitem{aviagen}
\bibitem{ross}
\bibitem{concept_model}
\bibitem{simulation_model}
\bibitem{model}
\bibitem{modelovani}
\bibitem{petriho_sit}
\bibitem{prechod}



\end{thebibliography}